\begin{resumo}[Abstract]
\begin{otherlanguage*}{english}

Factor models are widely used in asset pricing: they decompose cross-sectional returns into exposures to a reduced set of common variables, concentrating the market's relevant information without imposing prior economic structure. Recent advances in machine learning have allowed these factors to be extracted directly from data, treating them as latent variables learned jointly with the return prediction task. \textit{FactorVAE} extends this approach by modeling latent factors as random variables within a variational autoencoder, which allows the model to explicitly represent the uncertainty associated with the extracted signal, a relevant property given the low signal-to-noise ratio of financial time series. This work replicates and evaluates \textit{FactorVAE} in the Brazilian stock market, adapting the original framework to the B3 asset universe. The central question is whether the cross-sectional ranking advantage the model exhibits in the Chinese market is preserved in an environment with lower liquidity and fewer traded assets. Performance is assessed against three baselines (GRU, CA, IPCA), and results are analyzed in terms of both predictive ability and portfolio construction.

\textbf{Keywords}: \thekeywordsabstract.

\end{otherlanguage*}
\end{resumo}

\newpage